\documentclass[12pt]{article}
\usepackage{graphicx}
\usepackage[margin=2cm, a4paper]{geometry}
\usepackage{setspace}
\usepackage{pdfpages}
\usepackage{float}
\usepackage{amsmath}
\usepackage{fancyvrb}
\usepackage{amssymb}
\usepackage{minted}
\usepackage{enumitem}
\usepackage[colorlinks,linkcolor=blue]{hyperref}


\renewcommand{\contentsname}{Contents}
\renewcommand{\figurename}{Figure}
\renewcommand{\tablename}{Table}
\hypersetup{
    colorlinks=true,
    linkcolor=black,
    filecolor=magenta,      
    urlcolor=blue,
}
\newcommand{\mytitle}{Probability - Homework 1}
\newcommand{\myauthor}{B13902022 Yu-Chi,Lai}

\usepackage{fancyhdr}
\pagestyle{fancy}
\fancyhead{}
\fancyhead[L]{\mytitle}
\fancyhead[R]{\myauthor}

\title{\mytitle}
\author{\textbf{\myauthor}}
\date{}

\begin{document}

\maketitle

\section*{1.2-17}
\subsection*{(a)}
To draw 4 cards with the same face value, there's $13$ possibilities, then draw one from the rest of cards, there's $52-4=48$ possibilities (By definition, it's impossible to have same value with the previous $4$ cards.), Hence, the probability is:
\[\frac{13\times 48}{\binom{52}{4}}=\frac{1}{4165}\]
\subsection*{(b)}
We can choose two type of from $13$ face values, and for each face value we have $4$ colors to choose, then we can choose either $2$ or $3$ cards from them (but not for both types). (btw, we should time a extra $2$, for instance, if we need to construct a full-house using $2$ and $3$, we can draw $2$ cards from face value $2$ and $3$ cards from face value $3$, or draw $3$ cards from face value $2$ and $2$ cards from face value $3$, there's two possibilities.) Thus, the probability is:
\[\frac{\binom{13}{2}\binom{4}{2}\binom{4}{3}\times 2}{\binom{52}{5}}=\frac{6}{4165}\]
\subsection*{(c)}
Choose $3$ ranks from $13$ ($\binom{13}{3}$), then decide which of the $3$ ranks forms the triple ($\times 3$, since any of the $3$ chosen ranks can be assigned the triple role), choose $3$ suits for it ($\binom{4}{3}$), and choose $1$ suit for each of the remaining $2$ kicker ranks ($\binom{4}{1}^2$). Thus, the probability is:
\[\frac{\binom{13}{3}\binom{4}{3}(\binom{4}{1})^2\times 3}{\binom{52}{5}}=\frac{88}{4165}\]
\subsection*{(d)}
Choose $3$ ranks from $13$ ($\binom{13}{3}$), then decide which of the $3$ ranks is the single kicker ($\times 3$, since any of the $3$ chosen ranks can be assigned the kicker role), choose $1$ suit for it ($\binom{4}{1}$), and choose $2$ suits for each of the two pairs ($\binom{4}{2}^2$). Thus, the probability is:
\[\frac{\binom{13}{3}\binom{4}{1}(\binom{4}{2})^2\times 3}{\binom{52}{5}}=\frac{198}{4165}\]
\subsection*{(e)}
Choose $4$ ranks from $13$ ($\binom{13}{4}$), then decide which of the $4$ ranks forms the pair ($\times 4$, since any of the $4$ chosen ranks can be assigned the pair role), choose $2$ suits for it ($\binom{4}{2}$), and choose $1$ suit for each of the $3$ remaining kicker ranks ($\binom{4}{1}^3$). Thus, the probability is:
\[\frac{\binom{13}{4}\binom{4}{2}(\binom{4}{1})^3\times 4}{\binom{52}{5}}=\frac{352}{833}\]
\section*{1.3-13}
\subsection*{(a)}
\subsection*{(b)}
\subsection*{(c)}
\subsection*{(d)}
\subsection*{(e)}
\section*{1.4-15}
Define the random variable $X$ as: The fourth white balls appeaing at $X$-th round. And $P$ be the probability function.
\subsection*{(a)}
With replacement, the probability of taking white balls $p$ is $\frac{1}{2}$. Also we can consider the experiment as a bernoulli trial.
\begin{align*}
    P(X=4)=(1-p)^3p=\frac{1}{16} \\
    P(X=5)=(1-p)^4p=\frac{1}{32} \\
    P(X=6)=(1-p)^5p=\frac{1}{64} \\
    P(X=7)=(1-p)^6p=\frac{1}{128}
\end{align*} 
\subsection*{(b)}
We can consider the results as the permutations of $10$ white balls and $10$ red balls. And if we take $d$-th white ball at $r$-th round (where $d \le r$), the number of permutations is:
\[\frac{(r-1)!}{(d-1)!(r-d)!}\times\frac{(20-r)!}{(10-d)!(10-r+d)!}=\binom{r-1}{d-1}\times \binom{20-r}{10-d}\]
The total number of permutations of the $20$ balls is :
\[\frac{20!}{10!10!}=\binom{20}{10}\]

By the formula above, we can easily obtain all the answers:

\begin{align*}
P(X=4)= \\    
\end{align*}
\subsection*{(c)}
\section*{1.5-4}
Consider $4$ partitions (age 16-25, 26-50, 51-65, 66-90, respectively) as $B_1$, $B_2$, $B_3$, $B_4$, and $A$ be the event of accident. By Bayes' Theorem we obtain the answer:
\begin{align*}
    P(B_1|A) &= \frac{P(A|B_1)P(B_1)}{P(A|B_1)P(B_1)+P(A|B_2)P(B_2)+P(A|B_3)P(B_3)+P(A|B_4)P(B_4)} \\
    &= \frac{0.05\times 0.10}{0.05\times 0.10+0.02\times 0.55+0.03\times 0.20+0.04\times 0.15} \\
    &= 0.178571428571\approx 0.1786
\end{align*}


\end{document}
% how to display codes?
% \begin{minted}[frame=lines,framesep=2mm,baselinestretch=1.2,linenos,breaklines]{python}
% \end{minted}

% how to display images?
% \begin{figure}[H]
%     \centering
%     \includegraphics[width=0.5\linewidth]{}
%     \caption{Caption}
% \end{figure}
% test